%% =============================================================================
%% PAPER 1 - Section 6: Cross-Strategy Interdependencies
%% =============================================================================
%% Full prose. Tables/figures referenced via \ref{} only.
%% Actual tables in paper1_tables_sec6.tex; figures in paper1_figures_sec6.tex.
%% Maps to: rq3_01 through rq3_07
%% =============================================================================


%% ---- 6.1 Correlations, Co-Widening, and Persistence ----
\subsection{Correlations, Co-Widening, and Persistence}
\label{sec:rq3_corr}

% TODO: Full prose covering:
%
% Mispricing magnitude construction (rq3_01):
%   m_{i,t} = basis levels (not returns), Duffie (2010) framework
%   Negate CDS-bond basis so Delta_m > 0 = widening for all
%   BTP Italia: cross-sectional mean with maturity filter
%   CDS-Bond: top_n_negative
%   Regime definition: iTraxx Main fixed thresholds (60/120)
%
% Unconditional correlations (rq3_02 Section A):
%   Pearson + Spearman on returns, Delta_m, and levels m
%   HAC-robust inference (Newey-West, 6 lags)
%   Forbes-Rigobon bias correction for heteroskedasticity
%   → Table \ref{tab:rq3_unconditional_corr}
%
% Regime-dependent correlations (rq3_02 Section C):
%   Correlations by regime (LOW/NORMAL/HIGH)
%   Block bootstrap test for difference in co-widening across regimes
%   → Table \ref{tab:rq3_regime_corr}, Figure \ref{fig:rq3_regime_bar}
%
% Persistence analysis (rq3_02 Section D):
%   AR(1) coefficient phi per regime
%   Half-life of mispricing by regime
%   Duffie prediction: persistence higher in stress
%
% Rolling correlations (rq3_02 Section B):
%   12-month and 24-month windows
%   Visual evidence of time-varying co-movement


%% ---- 6.2 Spanning Regressions Across Strategies ----
\subsection{Spanning Regressions Across Strategies}
\label{sec:rq3_crossreg}

% TODO: Full prose covering:
%
% Fleckenstein-style spanning regressions (rq3_03 Sections A-B):
%   Delta_m_i = alpha + beta * Delta_m_j + epsilon
%   Contemporaneous + lagged specifications
%   HAC standard errors (Newey-West, 6 lags)
%   → Table \ref{tab:rq3_spanning}
%
% Interaction regressions (rq3_03 Section D):
%   Delta_m_i = alpha + beta_1*Delta_m_j + beta_2*(Delta_m_j * z_t) + beta_3*z_t + eps
%   z_t continuous (iTraxx Main standardized) + dummy HIGH
%   beta_2 > 0 implies stronger transmission during stress
%   → Table \ref{tab:rq3_interaction}
%
% PCA common factor structure (rq3_03 Section E):
%   PCA on Delta_m of the 3 strategies
%   PC1 variance share — common slow-moving capital factor?
%   PC1 regressed on funding stress proxies
%   → Table \ref{tab:rq3_pc1_funding}


%% ---- 6.3 VAR Analysis and Impulse Responses ----
\subsection{VAR Analysis and Impulse Responses}
\label{sec:rq3_var}

% TODO: Full prose covering:
%
% VAR on Delta_m: lag selection (BIC), estimation
% Granger Causality tests (rq3_04 Section B)
%   → Table \ref{tab:rq3_granger}
%
% GIRF — Generalized Impulse Response Functions (Pesaran-Shin 1998)
%   Invariant to ordering, bootstrap CI
%   → Figure \ref{fig:rq3_girf}
%
% Forecast Error Variance Decomposition (rq3_04 Section E)
%   Rising cross-strategy share at longer horizons
%   → Figure \ref{fig:rq3_fevd}
%
% Robustness in Appendix \ref{app:rq3_additional}:
%   All Cholesky permutations (rq3_04 Section D)
%   Threshold VAR: regime-dependent IRF (rq3_04 Section F)
%   VARX: exogenous stress proxy (rq3_04 Section G)


%% ---- 6.4 Discussion: Evidence for Slow-Moving Capital ----
\subsection{Discussion: Evidence for Slow-Moving Capital}
\label{sec:rq3_discussion}

% TODO: Full prose covering:
%
% Duffie scorecard — 6 testable predictions
%   P1. Positive cross-strategy correlation in Delta_m
%   P2. Co-widening in excess (P_joint > P_indep)
%   P3. Spanning R^2 > 0
%   P4. Granger causality (bidirectional)
%   P5. PC1 > 50% variance
%   P6. Persistence higher in stress (phi_HIGH > phi_LOW)
%   → Table \ref{tab:rq3_duffie_scorecard}, Figure \ref{fig:rq3_scorecard}
%
\paragraph{Scorecard Assessment.}
Table~\ref{tab:rq3_duffie_scorecard} summarises the evidence against six
testable predictions derived from the slow-moving capital framework of
\citet{duffie2010presidential}.  In the full sample, three of the six
predictions are confirmed at conventional significance levels.  Prediction~P1
(positive cross-strategy correlation in $\Delta m$) is confirmed for the
corporate credit cluster and, with opposite sign, for pairs involving
BTP~Italia, consistent with the two-cluster structure discussed below.
Prediction~P3 (spanning regressions with $R^{2}>0$) holds for three of the
six directional pairs, with particularly strong results for the
iTraxx$\,\to\,$BTP~Italia channel.  Prediction~P4 (Granger causality)
survives both Benjamini--Hochberg and Holm--Bonferroni corrections for
multiple testing, with the iTraxx$\,\to\,$BTP~Italia link emerging as the
single most robust result in the analysis.

The remaining predictions are limited by sample characteristics rather than
by contradictory evidence.  Prediction~P2 (excess co-widening in stress) is
confirmed for the corporate credit pair---co-widening rises by 14 to
22~percentage points in the \textsc{high} regime---but not for pairs
involving BTP~Italia, whose basis moves in the opposite direction during
stress.  Prediction~P5 (PC1 explaining more than 50\% of the variance in
$\Delta m$) is narrowly missed: PC1 accounts for 44.8\% of the variance,
and the 50\% threshold is conservative for a three-variable system.
Importantly, the economic content of PC1 is confirmed: it is significantly
explained by funding-stress variables (Table~\ref{tab:rq3_pc1_funding}),
establishing a direct link between the common factor driving correlated
mispricings and the intermediary balance-sheet conditions emphasised by
\citet{he2017intermediary}.  Prediction~P6 (higher persistence in stress)
cannot be reliably tested because only 8--9~months fall in the \textsc{high}
regime, which is insufficient for within-regime AR(1) estimation.  The
unconditional persistence, however, is substantial: half-lives range from
1.2~months (iTraxx) to 4.0~months (CDS--Bond), confirming that arbitrage
opportunities are slow to close.

Figure~\ref{fig:rq3_scorecard} shows that the scorecard is broadly stable
across sub-periods, with the corporate credit cluster consistently
satisfying more predictions than the full trivariate system.

\paragraph{Investor Base Segmentation: The Two-Cluster Structure.}
The central finding of this section is that the three strategies do not form a
single homogeneous group but two distinct clusters, and the key determinant is
the composition of the investor base.

The CDS--Bond basis and iTraxx Combined both operate within European corporate
credit, and their participants are overwhelmingly institutional: hedge funds
running relative value books, proprietary trading desks at dealer banks, and
leveraged accounts that rely on repo funding and derivatives margining.  These
investors share the same prime brokers, the same clearing infrastructure (ICE
Clear Credit, LCH), and the same balance sheet constraints.  When funding
conditions tighten---a margin call at a prime broker, a spike in repo haircuts,
or a redemption wave from fund investors---the same capital is withdrawn from
both markets simultaneously.  This is precisely the mechanism formalised by
\citet{BrunnermeierPedersen2009}: a negative funding shock forces leveraged
arbitrageurs to unwind positions across all markets where they are active,
generating correlated widening of mispricings.  Our finding that co-widening
between CDS--Bond and iTraxx rises by 14 to 22~percentage points in
\textsc{high} stress, and that the interaction coefficient amplifies by an order
of magnitude ($\hat{\beta}_{2} = -10.91$), is a direct empirical counterpart of
the margin-spiral prediction.

The BTP~Italia market is fundamentally different.  BTP~Italia bonds are
inflation-linked retail securities issued by the Italian Treasury specifically
to attract domestic household savings.  The \textit{Ministero dell'Economia e
delle Finanze} has consistently reported that approximately 65--70\% of
outstanding BTP~Italia are held by Italian retail investors, who purchase them
during dedicated placement windows attracted by the inflation protection and the
loyalty bonus (\textit{premio fedelt\`{a}}) paid at maturity to holders who
retain the bond from issuance.  These investors do not employ leverage, are not
subject to margin calls, do not use repo financing, and do not rebalance in
response to Value-at-Risk signals or institutional mandates.  They are, in the
language of \citet{VayanosVila2021}, ``preferred-habitat'' investors whose
demand is driven by personal savings decisions rather than by relative value
considerations.  As a consequence, the inflation-linked basis can persist
without attracting arbitrage capital, because the institutional relative-value
accounts that could close it represent only a small fraction of the outstanding
supply.  If anything, a flight-to-safety or a search for inflation protection
by retail investors may \emph{compress} the BTP~Italia basis precisely when
corporate credit bases are widening---generating the negative correlation
($\rho = -0.234$, $p = 0.018$) that we document.

This two-cluster result echoes the international evidence in
\citet{FontaineNolin2019}, who construct model-free indices of limits of
arbitrage across eight sovereign bond markets and find an analogous
segmentation: the indices of the United States, United Kingdom, Germany, Canada,
and Switzerland are strongly correlated with each other
($\rho = 0.44$--$0.86$), whereas Italy and France form a separate cluster
(mutually correlated at $\rho = 0.62$ but largely uncorrelated with the
Anglo-Saxon group).  The authors attribute this to the idiosyncratic impact of
the European sovereign debt crisis on Italian and French government bond
markets.  Our finding generalises this pattern: market segmentation in the
geography of slow-moving capital is determined not by the label of the asset
class but by the network of financial intermediaries---and ultimately by
\emph{who holds the bonds}.

\paragraph{Limits to Arbitrage in Corporate Credit.}
The persistence of the CDS--Bond basis is well understood in the literature as
a consequence of implementation frictions that prevent leveraged arbitrageurs
from fully closing the gap between cash bond and CDS spreads.
\citet{bai2019cds} provide the most detailed cross-sectional
decomposition: using Fama--MacBeth regressions on individual U.S.\ corporate
bonds from 2006 to 2014, they show that four limit-to-arbitrage channels
jointly explain approximately 37\% of the cross-sectional variation in the
negative basis at the bond level and roughly 77\% at the portfolio level during
the post-Lehman crisis.  The four channels are funding risk (the arbitrageur
must roll overnight-to-term repo financing at a haircut-dependent cost that
rises with balance-sheet stress), bond illiquidity (higher bid-ask spreads and
the transitory price-impact component of \citet{bao2011illiquidity} increase the
cost of entering and unwinding the cash leg), counterparty risk (the CDS
protection seller may default jointly with the reference entity, rendering the
hedge worthless at the worst possible moment), and collateral quality (lower
credit ratings imply larger haircuts and higher repo rates, reducing effective
leverage per unit of risk capital).  Perhaps most strikingly, bond lending fees
enter with a negative sign: rather than proxying for collateral value, high
lending fees during the crisis signal heavy short interest, and arbitrageurs
appear to avoid these bonds for fear that the basis will diverge further---a
``catching a falling knife'' effect that is itself a manifestation of
slow-moving capital.

\citet{oehmke2017anatomy} complement this cross-sectional evidence with a
quantity-based perspective.  Using net notional CDS positions from the DTCC,
they show that positions are significantly larger when the CDS--bond basis is
more negative, confirming that arbitrageurs lean against the mispricing.
Crucially, the sensitivity of CDS positions to the negative basis weakens when
funding conditions are tight, as measured by the TED spread: basis traders are
less active precisely when the opportunities are largest, because the same
funding shock that widens the basis also constrains the capital available to
exploit it.  When basis traders are active, they compress the negative basis and
reduce price impact in the bond market, acting as ``shock absorbers'' for bond
supply shocks.

This dynamic is exactly what \citet{mitchell2012arbitrage} document at the
aggregate level during the 2008 crisis: hedge funds running basis trades were
forced to liquidate when prime brokers withdrew financing, and the resulting
fire sales amplified the very dislocations the trades were designed to exploit.
The basis began to revert only after the U.S.\ government injected capital
through the Troubled Asset Relief Program (TARP), providing direct evidence
that the speed of convergence depends on the speed of capital replenishment.
Our European CDS--Bond basis strategy, while constructed on a different
universe of reference entities, operates through identical economic
mechanisms and is subject to the same funding-liquidity spirals.

\paragraph{Regulatory Capital and Index Replication Costs.}
The iTraxx index skew presents a third, and to our knowledge less studied,
channel through which limits to arbitrage sustain persistent mispricings in
European credit markets.  The CDS index is the instrument of first resort for
macro hedging: \citet{oehmke2017anatomy} document that CDS markets function
as ``alternative trading venues'' whose standardisation and liquidity advantages
attract both hedging and speculative flows, with speculative trading
concentrating almost exclusively in the CDS rather than the underlying bonds.
Index CDS contracts, as the most standardised instruments in the credit
derivative market, are the natural vehicle for this activity.  When investors
seek aggregate credit protection---whether to hedge portfolio exposures or to
express a directional view on credit spreads---they buy index CDS with little
regard for the fundamentals of individual constituent names.  This generates a
persistent demand pressure that pushes the index spread away from the intrinsic
value implied by the PV01-weighted average of the underlying 125~single-name
CDS.

Closing this gap requires assembling all 125~offsetting single-name positions
simultaneously, each with its own margin, counterparty, and clearing
arrangements.  This is operationally complex and capital-intensive.
\citet{boyarchenko2016trends} quantify the economics of this trade
from the dealer's perspective within a stylised balance-sheet framework.  They
show that post-crisis regulatory changes---in particular the Supplementary
Leverage Ratio (SLR) finalised in September 2014---have fundamentally altered
the cost structure of credit basis trades.  At a 6\% leverage ratio, which
corresponds to the effective SLR requirement for the largest U.S.\ bank
holding companies, the return on equity for a CDS-bond basis trade falls to
approximately 3\%, far below the 15\% internal hurdle rate that dealers
reportedly target.  Moreover, the breakeven CDS-bond basis required to achieve
a 15\% ROE shifts from $-134$ basis points under a 1\% leverage ratio to
$-218$ basis points under a 6\% ratio---a near-doubling of the dislocation
required before the trade becomes attractive.

The CDX-CDS skew trade faces an additional structural impediment that the
CDS-bond trade does not: index CDS contracts are subject to mandatory central
clearing, whereas single-name CDS contracts generally are not.  This asymmetry
in clearing mandates means that the margin requirements on the index leg
(posted with ICE) and the single-name legs (posted bilaterally or with a
different CCP) cannot be netted against each other.  The resulting capital
inefficiency further inflates the cost of carrying the trade.
\citet{boyarchenko2016trends} also document that investment advisors
have increasingly used index CDS contracts as substitutes for cash bond
purchases, selling protection on the index to gain credit exposure while
retaining liquidity in anticipation of fund outflows.  This strategic
positioning concentrates liquidity even more heavily in the index, draining
it from single-name CDS and widening the skew.

These regulatory and structural frictions explain why the iTraxx skew persists
even in benign market conditions and why it amplifies sharply during stress:
when funding conditions deteriorate, dealers not only refrain from new skew
arbitrage but may unwind existing positions, further widening the basis---the
same procyclical dynamic documented by \citet{mitchell2012arbitrage} for the
CDS-bond basis during the 2008 crisis.  Our finding that the iTraxx skew
belongs to the same cluster as the CDS--Bond basis, and that both widen
simultaneously in stress (with the interaction coefficient $\hat{\beta}_{2}$
amplifying by an order of magnitude in the \textsc{high} regime), is fully
consistent with this regulatory-capital channel operating across both segments
of the European corporate credit market.

\paragraph{From Alpha to Market Structure.}
The three limits to arbitrage identified above---retail investor segmentation,
funding and liquidity frictions, and regulatory capital costs---map directly
onto the alpha documented in Sections~4 and~5.  If capital were fast-moving and
unconstrained, the mispricings underlying all three strategies would be
arbitraged away rapidly, and the excess returns documented in our spanning and
conditional-alpha regressions would vanish.  The persistence of statistically
significant alpha, even after risk adjustment via rolling PCA
(Section~\ref{sec:rq1_pca}) and Adaptive Elastic Net
(Section~\ref{sec:rq2_aen}), is therefore not evidence of market inefficiency
in the classical sense but rather compensation for bearing the implementation
risks that prevent full arbitrage convergence.

This interpretation aligns with the broader theoretical literature on
intermediary-based asset pricing.  \citet{he2017intermediary} show that when
intermediaries are the marginal investors, asset risk premia depend on the
intermediary capital ratio rather than on consumption risk alone.
\citet{garleanu2011margin} demonstrate that margin constraints generate
deviations from the law of one price proportional to the difference in funding
requirements between otherwise identical assets.  And
\citet{mitchell2007slow} provide the original empirical evidence
that, following negative shocks to arbitrageur capital, prices of ``close
substitute'' assets diverge and then converge \emph{slowly}---at a speed
determined by the pace at which new risk capital enters the market, not by the
speed of information diffusion.

\citet{duffie2010presidential} synthesises these insights into the
slow-moving capital framework that motivates our RQ3.  The testable
implications he derives---correlated mispricings, predictable convergence,
cross-market transmission, and heightened co-movement during stress---are
precisely the patterns we document in Sections~6.1 through~6.3.  The three
strategies in our sample span what might be termed a ``liquidity pecking
order'' of increasing implementation cost: BTP~Italia requires no leverage and
no derivatives (buy and hold), the iTraxx skew is derivatives-only but requires
capital-intensive replication of 125~single-name positions, and the CDS--Bond
basis demands repo funding with rollover risk, CDS counterparty exposure, and
bond-market illiquidity.  The ranking of unconditional alpha across the three
strategies---BTP~Italia the lowest, CDS--Bond the highest---is qualitatively
consistent with this pecking order, suggesting that the cross-section of
arbitrage returns reflects a cross-section of implementation costs, exactly as
the slow-moving capital framework predicts.